\documentclass[10pt,a4paper]{amsart}
\usepackage{amsmath,amssymb,amsthm,graphicx}
% Retained minimal preamble for the Stacks Project chapter varieties.tex.
% Verbatim excerpt of the chapter preamble: Xy-pic setup (native preamble lines 14-18), the native theorem
% declarations with the mathematical macros (native preamble lines 154-217), and the project own no-op
% definitions of the reference/slogan/history display environments (native wrapper.tex lines 1-3) that the
% counted statement uses for its slogan. Copied mechanically; nothing added, removed or redefined.
% For commutative diagrams we use Xy-pic
\usepackage[all]{xy}

% We use 2cell for 2-commutative diagrams.
\xyoption{2cell}
% Theorem environments.
%
\theoremstyle{plain}
\newtheorem{theorem}[subsection]{Theorem}
\newtheorem{proposition}[subsection]{Proposition}
\newtheorem{lemma}[subsection]{Lemma}

\theoremstyle{definition}
\newtheorem{definition}[subsection]{Definition}
\newtheorem{example}[subsection]{Example}
\newtheorem{exercise}[subsection]{Exercise}
\newtheorem{situation}[subsection]{Situation}

\theoremstyle{remark}
\newtheorem{remark}[subsection]{Remark}
\newtheorem{remarks}[subsection]{Remarks}

\numberwithin{equation}{subsection}

% Macros
%
\def\lim{\mathop{\mathrm{lim}}\nolimits}
\def\colim{\mathop{\mathrm{colim}}\nolimits}
\def\Spec{\mathop{\mathrm{Spec}}}
\def\Hom{\mathop{\mathrm{Hom}}\nolimits}
\def\Ext{\mathop{\mathrm{Ext}}\nolimits}
\def\SheafHom{\mathop{\mathcal{H}\!\mathit{om}}\nolimits}
\def\SheafExt{\mathop{\mathcal{E}\!\mathit{xt}}\nolimits}
\def\Sch{\mathit{Sch}}
\def\Mor{\mathop{\mathrm{Mor}}\nolimits}
\def\Ob{\mathop{\mathrm{Ob}}\nolimits}
\def\Sh{\mathop{\mathit{Sh}}\nolimits}
\def\NL{\mathop{N\!L}\nolimits}
\def\CH{\mathop{\mathrm{CH}}\nolimits}
\def\proetale{{pro\text{-}\acute{e}tale}}
\def\etale{{\acute{e}tale}}
\def\QCoh{\mathit{QCoh}}
\def\Ker{\mathop{\mathrm{Ker}}}
\def\Im{\mathop{\mathrm{Im}}}
\def\Coker{\mathop{\mathrm{Coker}}}
\def\Coim{\mathop{\mathrm{Coim}}}

% Boxtimes
%
\DeclareMathSymbol{\boxtimes}{\mathbin}{AMSa}{"02}

%
% Macros for moduli stacks/spaces
%
\def\QCohstack{\mathcal{QC}\!\mathit{oh}}
\def\Cohstack{\mathcal{C}\!\mathit{oh}}
\def\Spacesstack{\mathcal{S}\!\mathit{paces}}
\def\Quotfunctor{\mathrm{Quot}}
\def\Hilbfunctor{\mathrm{Hilb}}
\def\Curvesstack{\mathcal{C}\!\mathit{urves}}
\def\Polarizedstack{\mathcal{P}\!\mathit{olarized}}
\def\Complexesstack{\mathcal{C}\!\mathit{omplexes}}
% \Pic is the operator that assigns to X its picard group, usage \Pic(X)
% \Picardstack_{X/B} denotes the Picard stack of X over B
% \Picardfunctor_{X/B} denotes the Picard functor of X over B
\def\Pic{\mathop{\mathrm{Pic}}\nolimits}
\def\Picardstack{\mathcal{P}\!\mathit{ic}}
\def\Picardfunctor{\mathrm{Pic}}
\def\Deformationcategory{\mathcal{D}\!\mathit{ef}}
\newenvironment{reference}{}{}
\newenvironment{slogan}{}{}
\newenvironment{history}{}{}
\usepackage{cleveref}
\crefname{originaldefinition}{original definition}{original definitions}
\Crefname{originaldefinition}{Original definition}{Original definitions}
\makeatletter
\newlabel{morphisms-definition-dominant-rational}{{0A1Z}{source}{}{sourcepointer-1}{}}
\newlabel{morphisms-definition-dominant-rational@cref}{{[originaldefinition][0][]0A1Z}{[1][0][]source}{}{}{}}
\makeatother
\crefname{originallemma}{original lemma}{original lemmas}
\Crefname{originallemma}{Original lemma}{Original lemmas}
\makeatletter
\newlabel{morphisms-lemma-rational-map-finite-presentation}{{0BX8}{source}{}{sourcepointer-2}{}}
\newlabel{morphisms-lemma-rational-map-finite-presentation@cref}{{[originallemma][0][]0BX8}{[1][0][]source}{}{}{}}
\makeatother
\title{Varieties with dominant rational maps are equivalent to finitely generated field extensions}
\author{The Stacks Project Authors}
\date{stacks-project:varieties.tex:local-snapshot-20261011; GNU FDL 1.2 or later}
\begin{document}
\maketitle
\noindent\textit{Editorial note.} All delivered mathematical text is the author's own, copied from the pinned Stacks Project chapter file varieties.tex; no statement or proof step was rewritten, reordered or removed.\par\medskip
\noindent\textit{Editorial note.} The preamble is the chapter's own retained minimal preamble (sources/stacks-varieties/retained-preamble.tex), a verbatim excerpt of the native Xy-pic setup, of the native theorem declarations with mathematical macros, and of the project's own no-op definitions of the reference display environments; no macro was invented or redefined.\par\medskip
\noindent\textit{Editorial note.} Both referenced results live in earlier chapters, so their locators point at the verbatim lines of the project's own tag list staged in the snapshot (sources/stacks-chow/tags.txt) that carry the canonical tags.\par\medskip
\noindent\textit{Editorial note.} Overlap pre-check: this label occurs in no other staged chapter (run with agent-authoring/stk/card\_plan.py varieties before opening).\par\medskip
\noindent\textit{Allowed standard prerequisites (editorial).}\begin{itemize}
\item Dominant rational map between varieties. Morphisms, Definition labelled morphisms-definition-dominant-rational (canonical Stacks tag 0A1Z) together with the discussion following it; invoked at proof lines 173-176. Supplies the characterisation of dominant rational maps as those rational maps sending the generic point of the source to the generic point of the target, which is what makes the functor on function fields well defined.
\item A local map of function fields determines a unique dominant rational map. Morphisms, Lemma labelled morphisms-lemma-rational-map-finite-presentation (canonical Stacks tag 0BX8); invoked at proof lines 184-186. Supplies the converse direction, hence full faithfulness of the functor.
\end{itemize}\medskip
\ifdefined\hypertarget\hypertarget{sourcepointer-1}{}\fi
\noindent\textit{Source reference.} Original definition 0A1Z: sources/stacks-chow/tags.txt, source lines 12687--12687. Cross-chapter pointer invoked at proof lines 175-176 for the characterisation of dominant rational maps at the generic point; the functor and the essential-image argument are retained in the counted span.\par\medskip
\ifdefined\hypertarget\hypertarget{sourcepointer-2}{}\fi
\noindent\textit{Source reference.} Original lemma 0BX8: sources/stacks-chow/tags.txt, source lines 14971--14971. Cross-chapter standard prerequisite invoked once at proof line 186 for the converse direction; the full-faithfulness conclusion drawn from it is retained in the counted span.\par\medskip
\begin{theorem}
\label{theorem-varieties-rational-maps}
Let $k$ be a field. The category of varieties and
dominant rational maps is equivalent to the category of
finitely generated field extensions $K/k$.
\end{theorem}

\begin{proof}
Let $X$ and $Y$ be varieties with generic points $x \in X$ and $y \in Y$.
Recall that dominant rational maps from $X$ to $Y$ are exactly those
rational maps which map $x$ to $y$
(Morphisms, Definition \ref{morphisms-definition-dominant-rational}
and discussion following).
Thus given a dominant rational map $X \supset U \to Y$ we obtain a map of
function fields
$$
k(Y) = \kappa(y) = \mathcal{O}_{Y, y}
\longrightarrow
\mathcal{O}_{X, x} = \kappa(x) = k(X)
$$
Conversely, such a $k$-algebra map (which is automatically local as the
source and target are fields) determines (uniquely) a dominant rational
map by Morphisms, Lemma \ref{morphisms-lemma-rational-map-finite-presentation}.
In this way we obtain a fully faithful functor.
To finish the proof it suffices to show that every finitely generated
field extension $K/k$ is in the essential image.
Since $K/k$ is finitely generated, there exists a finite type
$k$-algebra $A \subset K$ such that $K$ is the fraction field of $A$.
Then $X = \Spec(A)$ is a variety whose function field is $K$.
\end{proof}

\end{document}
